% Chapter opener (rewritten 2026-10-31, Aida-style single paragraph; refactored 2026-11-06 to nine
% stages, one per section, in \input order). Moved from thesis/methodology/00_chapter_intro.tex
% on 2026-11-07. \input order is in content/methods.tex:
%   01a-02b   full text
%   03-07     notes only, to be written
% Overview figure still to make (style guide: an overview figure opens the methods).
\providecommand{\clarify}[1]{\textcolor{red}{\textbf{Clarification needed:} #1}}
\chapter{Methodology}
\label{ch:methodology}

The methodology followed in this thesis is summarized in Figure~\ref{fig:method-overview} and consists of nine main stages:
(1) extraction of the coronary centerlines and lumen contours from all baseline scans with an in-house tracer;
(2) comparison of the extracted centerlines and lumen with Medis as a reference, as no ground truth was available;
(3) selection of the plaque-free baseline cohort and definition of the outcome from the visual analysis;
(4) division of the traced arteries into the segments of the visual analysis;
(5) the prediction model and its cross-validation;
(6) the number of parameters the cohort can support and the choice of geometric features within it;
(7) definition of the geometric features;
(8) assessment of their measurement properties, against synthetic shapes and known groups, and of the variance explained by known risk factors;
and (9) a secondary, descriptive analysis of the shape of the curvature profile.

\clarify{overview figure of the nine stages (label fig:method-overview).}
